% !TeX root = ../papers/nc-p3.tex
\section{直線ファイバーの中心は曲線以下である}
\label{sec:branching-curve-bound}
\begin{theorem}\label{thm:branching-centers-curve}
標準正則四窓の各辺、各射影について、直線ファイバーを持つ点の軌跡は次元1以下である。
従って一つの辺がある曲面の全点の上に直線を持つことはない。
純関係が一つもないという仮定は不要である。
\end{theorem}
\begin{proof}[一般平面における六関係の独立性]
$R\subset V\otimes W$ を辺の六関係とする。
どの $a\ne0$ に対しても $ka\otimes W$ に入る関係は高々二次元である。
三次元なら定理\ref{thm:all-endpoint-ranks-two}の証明で $k=3$ とした議論がそのまま指数成長を与え、
四次元なら右socleがある。逆の $V\otimes kb$ でも同じである。

一般の $a\in V$ については $R\cap(ka\otimes W)=0$ である。
反対なら $\mathbb P(R)$ 内の純関係の既約族の一つが $\mathbb P V$ に支配的となる。
第二の点 $b$ が一定なら $V\otimes kb\subset R$ で矛盾する。従ってその族の関数体で
$b$ の二座標の比 $f$ は非定数である。一座標を自明化した $a$ の四一次式の空間を $A$ とすると、
$\dim(A+fA)\ge7$ となる。実際、$f v=u$ と $f v'=u'$ から $u v'=u' v$ が
$\mathbb P V$ の斉次二次式の恒等式となる。$f$ が非定数なら $u,v$ は独立で互いに素なので
二対は比例する。よって $A\cap fA$ の次元は高々1である。
これら七つ以上の独立関数は $a\otimes b$ の座標の包に入り、六次元性に反する。
従って一般超平面 $a^\perp$ に制限すれば六関係は一次行列の列として独立のままである。

\end{proof}
\begin{proof}[二次・三次の曲面]
$r(p)=\dim R(p)$ と書く。$r(p)=2$ の被約軌跡上の像は階数2の部分束
$S\subset W\otimes\mathcal O$ であり、$S(1)$ は六関係で大域生成される。
$S^*$ の行列式の与えるPlücker写像は閉点上で単射である（定理\ref{thm:all-endpoint-ranks-two}）。
まず全三次小行列式が恒等零の場合は下の平面の場合に含める。
そうでなければ曲面成分は三次以下である。既約な二次または三次の曲面の一般平面切断を $C$ とすれば、
$C$ は整な平面曲線で、二次なら $\mathbb P^1$、三次なら特異な場合も算術種数1で $\omega_C=\mathcal O_C$ である。
六つの生成切断は $H^0(C,S(1))$ で独立である。$C$ は切断平面を張るから、
関係の一次行列が $C$ で消えればその平面で消え、前段が適用できる。
大域生成された階数2束の一般切断は曲線上で零点を持たない
（各点で二条件の入射対応の次元を数える）。従って
$h^0(S(1))\le1+h^0(\det S(1))$ である。$D=\det S^*$ とおけば
$\det(S(1))=\mathcal O_C(2)D^{-1}$ である。
二次なら $\deg D\ge1$ なので右辺は高々5となる。
整な平面三次なら $D$ は正次数であり、Riemann--Rochと双対性から $h^0(D)=\deg D$ である。
$\deg D=1$ は非定数性に反し、$\deg D=2$ は直線への二枚被覆で標数零の単射性に反する。
よって $\deg D\ge3$、$\mathcal O_C(2)D^{-1}$ の大域切断は高々3である。
いずれも六つの独立切断に足りない。
\end{proof}

\begin{proof}[平面の排除]
$r(p)=2$ が平面 $P=\mathbb P(a^\perp)$ 全体で成立すると仮定する。
$\det S^*=\mathcal O_P(e)$ とおく。$\det S(1)$ も大域生成され、単射なGrassmann写像より
$1\le e\le2$ である。$e=2$ では $S(1)$ は大域生成で行列式自明なので自明束である
（一点で独立な二切断の行列式は消えない定数）。
従って六関係のうち少なくとも四次元が $ka\otimes W$ に入るため不可能である。

$e=1$ ならPlücker写像は線型平面を $\operatorname{Gr}(2,W)$ に埋め込む。
そのような平面は全ての二空間が一定の三空間に入るものか、一定の直線を含むものである。
実際、分解可能な二ベクトルの一次結合が分解可能なら二直線は射影空間で交わる。
二本を固定すると、それらと交わる残る直線は同一平面に入るか、共通点を通る。既約性から二型を得る。
前者なら $S$ は $3\mathcal O\to\mathcal O(1)$ の核で、$h^0(S(1))=3$ である
（三個の独立な一次式の一次関係を数える）。
少なくとも三関係が $ka\otimes W$ に入り先の上界に反する。

後者の共通直線を $kb$ とする。$S$ は $\mathcal O$ と $\mathcal O(-1)$ の拡大で
$H^1(\mathcal O(1))=0$ から分裂し、$h^0(S(1))=4$ である。
$R$ の制限は四次元全部でなければ $ka\otimes W$ の関係が三つ以上になる。
従って $R$ 自体は非零の $a\otimes v$ を含み、二文字を
$V/ka,W/kb$ へ落とした後の $R$ の像は非退化な一つの双一次式 $\phi$ で張られる。
実際、$S/kb$ の与える $P\to\mathbb P(W/kb)$ は射影線型同型で、
対応する一次行列以外は全て $b$ に値を持つ。位置を順に0,1とすれば、両文字を
このように落とした四テンソルは $\phi\otimes C$、$0\ne C\in R_2$ である。
他の隣接切断ではこの制限の階数は $3\operatorname{rank} C\le6$ であり、$C$ の左支持 $U$ は高々二次元である。
任意の $y\in b^\perp\setminus0$, $z\in U^\perp$ で $w(y,z)\in R_3$ は
$V_0$ を $ka$ で割ると零なので $V_3\otimes ka$ に入る。
しかし $R_0$ の純関係 $a\otimes v$ と連続する純関係はない。従って $w(y,z)=0$ である。
ゆえに全ての $y\in\mathbb P(b^\perp)$ において $R_1(y)\subset U$ である。
その次元は少なくとも2なので $U$ に等しく、相異なる点は同じ端支持平面を持たないという定理に反する。
これで平面も全射影空間の場合も排除した。
\end{proof}

\begin{corollary}\label{cor:branching-curves-successor-surface}
直線ファイバーの中心が既約曲線をなす場合、それらの直線の合併の閉包は既約曲面である。
曲線の一般点では逆方向にも少なくとも一点が続く。前点を持たず直線の後続を持つ点は各位置に有限個しかない。
\end{corollary}
\begin{proof}
直線は中心から一意に決まり相異なるので、その射影直線束の像は二次元である。
曲線の一般点で逆方向の評価階数が4なら、辺の隣接切断の直交不等式により全ての
後続で順方向の評価階数は2以下となる。すると上の曲面全体が新しい直線ファイバーの中心に含まれて矛盾する。
逆方向の階数4は開条件で、各曲線に稠密にはなれないので、そのような点は有限個である。
\end{proof}
